%!TeX root = ../principal.tex
%!TeX encoding = utf-8
% =====================================================================
%  PROTOCOLO EXPERIMENTAL
%  Linhas marcadas com "% VERIFICAR" ou "% PREENCHER" pedem conferência
%  ou um valor que só existe nos arquivos locais (paper_info/REPORT.md).
% =====================================================================

\chapter{Protocolo Experimental}
\label{ch:protocol}

Este capítulo especifica o delineamento do experimento. Comparamos
\textsc{Efficient-CapsNet}, \textsc{ResNet-18} e \textsc{DeiT-Tiny} sob variação
controlada de três fatores: tamanho do conjunto de treino, intensidade da
augmentação e semente aleatória. O capítulo fecha com o procedimento de
avaliação sob transformações nunca vistas no treino.

% ---------------------------------------------------------------------
\section{Objetivos e hipóteses}
\label{sec:protocol:hypotheses}

Três perguntas orientam o experimento.

\begin{description}
\item[RQ1 (Equivariância).]
As cápsulas degradam menos que os modelos convolucionais e de atenção sob
transformações ausentes do treino, em particular sob grandes rotações?
\item[RQ2 (Eficiência de dados).]
A diferença de acurácia entre as arquiteturas cresce a favor das cápsulas
conforme a fração de dados de treino diminui?
\item[RQ3 (Interação com augmentação).]
Uma eventual vantagem das cápsulas é arquitetural, ou as linhas de base a
recuperam quando recebem augmentação explícita?
\end{description}

\noindent
Cada pergunta vira uma hipótese falsificável. A Seção~\ref{sec:protocol:stats}
define o teste que aceita ou rejeita cada uma.

\begin{description}
  \item[H1.] Sob grande rotação, a razão de retenção $R_c$ é maior para o
             Efficient-CapsNet que para as linhas de base, em comparações
             pareadas por semente.
  \item[H2.] A queda de acurácia de $f = 1{,}00$ para $f = 0{,}33$ é menor para
             o Efficient-CapsNet que para as linhas de base.
  \item[H3.] Os efeitos de H1 e H2 encolhem até o nível do ruído entre sementes
             quando o regime de augmentação forte é aplicado.
\end{description}

Uma decisão da análise afeta a leitura de H1. A
Seção~\ref{sec:res:confound} mostra que, nesta grade, $R_c$ é
anticorrelacionada com a acurácia limpa e favorece modelos fracos. Por isso H1
é julgada pela retenção pareada por semente \emph{e} pela acurácia absoluta sob
rotação, e não pela média de $R_c$ isoladamente.

% ---------------------------------------------------------------------
\section{Conjuntos de dados}
\label{sec:protocol:data}

\subsection{CIFAR-10}
O CIFAR-10~\citar{krizhevsky2009cifar} reúne 60.000 imagens coloridas de
\(32\times32\) pixels em 10 classes, 6.000 por classe. A divisão oficial separa
50.000 imagens para treino e 10.000 para teste, com 1.000 de cada classe no
teste.

As dez classes são avião, automóvel, pássaro, gato, cervo, cachorro, sapo,
cavalo, navio e caminhão.

Fonte: CIFAR-10, University of Toronto, acesso em 14 de agosto de 2026.

\subsection{D-Fire}
\label{sec:protocol:dfire}
O D-Fire~\citar{venancio2022dfire} é um conjunto de imagens para detecção de
fogo e fumaça, construído pelo grupo Gaia em colaboração com a Universidade
Federal de Minas Gerais. São 21.527 imagens de câmeras de vigilância, de
veículos aéreos não tripulados e da web, em cenas urbanas, industriais e
florestais.

Essa origem heterogênea motiva a escolha. Conjuntos capturados em uma única
campanha de aquisição variam pouco; o D-Fire varia em ponto de vista, escala e
domínio, que é o que interessa a um estudo sobre robustez.

As anotações originais seguem o formato YOLO, com caixas normalizadas para duas
categorias: \emph{fire} (14.692 caixas) e \emph{smoke} (11.865 caixas). A
Tabela~\ref{tab:dfire-content} mostra a distribuição das imagens por conteúdo.
Neste trabalho descartamos as coordenadas das caixas e usamos apenas duas
classes: \textsc{fire/smoke}, para imagens com ao menos uma caixa, e
\textsc{none}, para as demais.

\begin{table}[t]
  \centering
  \caption{Composição do D-Fire completo, como distribuído pelos autores, por
  conteúdo anotado, e o agrupamento binário adotado neste trabalho.}
  \label{tab:dfire-content}
  \begin{tabular}{lrl}
    \toprule
    Conteúdo anotado & Imagens & Classe binária \\
    \midrule
    Somente fogo               & 1.164  & \textsc{fire/smoke} \\
    Somente fumaça             & 5.867  & \textsc{fire/smoke} \\
    Fogo e fumaça              & 4.658  & \textsc{fire/smoke} \\
    Nenhum (\emph{background}) & 9.838  & \textsc{none} \\
    \midrule
    Total                      & 21.527 & \\
    \bottomrule
  \end{tabular}
\end{table}

Fonte: D-Fire Dataset, Gaia --- \url{https://github.com/gaiasd/DFireDataset};
cópia utilizada: \emph{Smoke-Fire-Detection-YOLO},
\url{https://www.kaggle.com/datasets/sayedgamal99/smoke-fire-detection-yolo},
acesso em 9 de setembro de 2026.

\subsection{Partições e pré-processamento}
\label{sec:protocol:splits}
Em ambos os conjuntos, uma fração fixa $\rho_{\text{val}} = 0{,}1$ da partição
de treino é separada aleatoriamente para validação e seleção do modelo. O
sorteio depende apenas da semente, e não da arquitetura nem da fração de dados:
as três arquiteturas veem a mesma validação para uma mesma semente. A fração de
dados $f$ é aplicada depois, sobre o que resta do treino, de modo que a
validação é idêntica entre as frações. O conjunto de teste serve apenas à
avaliação final do Capítulo~\ref{ch:results}.

No CIFAR-10 isso deixa 45.000 imagens de treino e 5.000 de validação, e o
treino efetivo tem 14.850, 29.700 e 45.000 imagens para $f = 0{,}33$, $0{,}66$
e $1{,}00$. As imagens mantêm a resolução nativa.

A cópia do D-Fire utilizada traz 14.122 imagens na partição de treino e 4.306 na
de teste; o treino original dos autores tem 17.221 imagens, e as demais estão
em uma pasta de validação da cópia que não foi usada.
% VERIFICAR: confirmar em paper_info/dfire_split_counts.csv que as ~3.099
% imagens restantes estão na pasta val/valid da cópia do Kaggle.
Após a separação da validação ficam 12.710 imagens de treino e 1.412 de
validação, e o treino efetivo tem 4.194, 8.388 e 12.710 imagens para as três
frações. Todas as imagens são redimensionadas para $224\times224$ pixels, sem
recorte, para não remover a região de fogo ou fumaça. Na partição de teste, a
classe majoritária responde por $53{,}4\%$ das imagens.

A validação do D-Fire passa apenas por redimensionamento e normalização. No
CIFAR-10, as imagens de validação passam pelo mesmo \emph{pipeline} de
augmentação do treino, de modo que a perda de validação usada na seleção do
modelo é medida sobre imagens aumentadas no regime de cada execução
(Seção~\ref{sec:protocol:threats}).

\begin{table}[t]
  \centering
  \caption{Conjuntos de dados utilizados. ``Treino'' é o que resta após a
  separação da validação, com $f = 1{,}00$.}
  \label{tab:datasets}
  \begin{tabular}{lccccc}
    \toprule
    Conjunto & Classes & Treino & Val. & Teste & Entrada \\
    \midrule
    CIFAR-10 & 10 & 45.000 & 5.000 & 10.000 & $3\times32\times32$ \\
    D-Fire   & 2  & 12.710 & 1.412 & 4.306  & $3\times224\times224$ \\
    \bottomrule
  \end{tabular}
\end{table}

\subsection{Treinamento do zero e capacidade}
Todos os modelos partem de inicialização aleatória, sem pré-treino em ImageNet
ou em qualquer outro conjunto. Pré-treino confundiria a comparação de
eficiência de dados, porque o modelo chegaria ao experimento com um volume de
dados que as outras arquiteturas não viram.

Não pareamos as arquiteturas por número de parâmetros nem por custo de
inferência. Cada uma é usada na configuração padrão de sua família, e a
Tabela~\ref{tab:model-complexity} registra o desnível que resulta. A ResNet-18
tem cerca de $2{,}8$ a $2{,}9$ vezes os parâmetros do Efficient-CapsNet, e o
DeiT-Tiny cerca de $1{,}4$ vez. O desnível é maior do que a tabela sugere,
porque cerca de $3{,}7$~M dos parâmetros capsulares estão no decodificador de
reconstrução (Seção~\ref{sec:protocol:models:ecaps}); o extrator e as cápsulas
somam apenas cerca de $0{,}16$~M.

\begin{table}[t]
  \centering
  \caption{Complexidade dos modelos. Parâmetros e operações de multiplicação e
  acumulação (MACs), contadas com a biblioteca \texttt{ptflops}~\citar{ptflops}
  em uma passagem direta com lote unitário, na resolução de entrada de cada
  conjunto. Para o Efficient-CapsNet a contagem inclui o decodificador, que
  também é executado na inferência.}
  \label{tab:model-complexity}
  \begin{tabular}{lcccc}
    \toprule
    \multirow{2}{*}{Modelo} & \multicolumn{2}{c}{CIFAR-10} & \multicolumn{2}{c}{D-Fire} \\
    \cmidrule(lr){2-3}\cmidrule(lr){4-5}
     & Params (M) & GMACs & Params (M) & GMACs \\
    \midrule
    Efficient-CapsNet & 3,94 & 0,049 & 3,85 & 0,458 \\
    ResNet-18         & 11,17 & 0,558 & 11,18 & 1,825 \\
    DeiT-Tiny         & 5,36 & 0,252 & 5,52 & 0,913 \\
    \bottomrule
  \end{tabular}
\end{table}

% ---------------------------------------------------------------------
\section{Protocolo de treinamento}
\label{sec:protocol:training}

\subsection{Função objetivo}
O modelo de cápsulas é treinado com perda de margem:
\begin{equation}
  \mathcal{L}_{k} = T_{k}\max(0, m^{+} - \lVert \mathbf{v}_{k}\rVert)^{2}
  + \lambda\,(1 - T_{k})\max(0, \lVert \mathbf{v}_{k}\rVert - m^{-})^{2},
  \label{eq:margin}
\end{equation}
com $m^{+} = 0{,}9$, $m^{-} = 0{,}1$ e $\lambda = 0{,}5$. Soma-se a ela um termo
de reconstrução quadrático médio, com peso $\alpha = 5\times10^{-4}$.

A ResNet-18 e o DeiT-Tiny são treinados com entropia cruzada. O objetivo
diferente é um confundidor da comparação arquitetural. Ele seria evitável
treinando as linhas de base também com perda de margem, variante que não foi
executada; o confundidor é retomado na Seção~\ref{sec:protocol:threats}.

\subsection{Otimização}
Todos os modelos compartilham a configuração da
Tabela~\ref{tab:hyperparams}, herdada da implementação de referência do
Efficient-CapsNet e usada sem busca de hiperparâmetros.
% VERIFICAR: confirmar que não houve ajuste de hiperparâmetros.
O escalonamento reduz a taxa de aprendizado a cada época, até cerca de
$6{,}8\times10^{-5}$ na última. A norma $\ell_2$ do gradiente é limitada a
$5{,}0$ antes de cada passo, e um passo com perda ou gradiente não finito
interrompe o treino. A parada antecipada está desativada: todo treino roda as
100 épocas, e o \emph{checkpoint} de menor perda de validação é o avaliado.

\begin{table}[t]
  \centering
  \caption{Hiperparâmetros de treino. Idênticos para todas as arquiteturas,
  frações de dados, regimes de augmentação e sementes.}
  \label{tab:hyperparams}
  \begin{tabular}{ll}
    \toprule
    Item & Valor \\
    \midrule
    Otimizador               & Adam \\
    Taxa de aprendizado      & $5\times10^{-4}$ \\
    Escalonamento            & Exponencial, $\gamma = 0{,}98$ por época \\
    Tamanho do lote          & 16 \\
    Épocas                   & 100 (sem parada antecipada) \\
    Decaimento de peso       & 0 (nenhum) \\
    Recorte do gradiente     & norma $\ell_2$ máx. 5,0 \\
    Seleção do modelo        & Menor perda de validação \\
    Precisão                 & fp32 (sem AMP) \\
    Fração de validação      & 0,1 \\
    \bottomrule
  \end{tabular}
\end{table}

\subsection{Orçamento de comparação}
\label{sec:protocol:budget}
Todas as comparações deste trabalho são feitas com o mesmo número de épocas.
Épocas iguais não significam custo igual: uma arquitetura mais lenta por época
consome mais tempo para o mesmo número de passos. Por isso o custo de cada
treino é reportado à parte, na Tabela~\ref{tab:cost}, e o leitor pode julgar
se uma diferença de acurácia compensa uma diferença de tempo. Uma comparação
com tempo de treino igual exigiria avaliar \emph{checkpoints} intermediários,
e não foi realizada.

\subsection{Hardware e tempo de execução}
Os modelos foram treinados em uma RTX 2060 de 6\,GB de VRAM. O ambiente usa
CUDA 13.3, PyTorch 2.13.0, Python 3.14.7 e Arch Linux (kernel 7.1.8-arch1-3).
As arquiteturas vêm da biblioteca \texttt{timm}~\citar{wightman2019timm}, com
exceção do Efficient-CapsNet, implementado neste trabalho.

% ---------------------------------------------------------------------
\section{Delineamento experimental}
\label{sec:protocol:design}

O estudo é um fatorial completo sobre cinco fatores, resumidos na
Tabela~\ref{tab:design}.

\begin{table}[t]
  \centering
  \caption{Fatores e níveis. A grade completa soma
  $3 \times 2 \times 3 \times 3 \times 3 = 162$ treinos.}
  \label{tab:design}
  \begin{tabular}{lll}
    \toprule
    Fator & Níveis & $|{\cdot}|$ \\
    \midrule
    Arquitetura             & Efficient-CapsNet, ResNet-18, DeiT-Tiny & 3 \\
    Conjunto de dados       & CIFAR-10, D-Fire                        & 2 \\
    Fração de dados $f$     & $0{,}33$, $0{,}66$, $1{,}00$            & 3 \\
    Augmentação $a$         & nenhuma, padrão, forte                  & 3 \\
    Semente $s$             & $1$, $2$, $3$                           & 3 \\
    \bottomrule
  \end{tabular}
\end{table}

\subsection{Ablação da augmentação}
\label{sec:protocol:aug}

A Tabela~\ref{tab:aug-regimes} lista as operações de cada regime. Os dois
conjuntos diferem em dois pontos. O D-Fire não usa recorte aleatório, que
poderia remover a única região com fogo ou fumaça e rotular a imagem errado. E
o regime forte do D-Fire acrescenta variação de brilho e contraste antes do
RandAugment~\citar{cubuk2020randaugment}.

\begin{table}[t]
  \centering
  \caption{Regimes de augmentação em treino. Todos são seguidos de conversão
  para tensor e normalização por canal; no D-Fire, todos são precedidos de
  redimensionamento para $224\times224$. O Random Erasing é aplicado depois da
  normalização.}
  \label{tab:aug-regimes}
  \small
  \begin{tabular}{l l p{0.62\textwidth}}
    \toprule
    Regime & Conjunto & Operações \\
    \midrule
    nenhuma & ambos & --- \\
    \midrule
    \multirow{2}{*}{padrão}
      & CIFAR-10 & recorte aleatório $32\times32$ (preenchimento 4),
                   espelhamento horizontal ($p=0{,}5$), rotação $\pm15^{\circ}$ \\
      & D-Fire   & espelhamento horizontal ($p=0{,}5$), rotação $\pm15^{\circ}$ \\
    \midrule
    \multirow{2}{*}{forte}
      & CIFAR-10 & padrão $+$ RandAugment ($N=2$ operações, magnitude $M=9$ de
                   30) $+$ Random Erasing~\citar{zhong2020randomerasing}
                   ($p=0{,}5$, área $0{,}02$--$0{,}33$, razão
                   $0{,}3$--$3{,}3$, valor 0) \\
      & D-Fire   & padrão $+$ brilho e contraste $\times[0{,}7;\,1{,}3]$ $+$
                   RandAugment ($N=2$, $M=9$) $+$ Random Erasing (mesmos
                   parâmetros) \\
    \bottomrule
  \end{tabular}
\end{table}

A rotação de treino vai a $\pm15^{\circ}$ no regime padrão e a cerca de
$\pm24^{\circ}$ no forte, somando a rotação do RandAugment ($9^{\circ}$ na
magnitude usada). A sonda de rotação da Seção~\ref{sec:protocol:rotation}
sorteia ângulos até $\pm90^{\circ}$, de modo que a maior parte das amostras cai
fora do que o treino cobre, mas não todas: cerca de $17\%$ dos ângulos
sorteados estão dentro da faixa do regime padrão e cerca de $27\%$ dentro da do
regime forte.

% ---------------------------------------------------------------------
\section{Protocolo de avaliação}
\label{sec:protocol:eval}

Cada modelo é avaliado no \emph{checkpoint} de menor perda de validação, sob três
condições.

\subsection{Avaliação em distribuição}
O teste limpo passa apenas por normalização (e, no D-Fire, redimensionamento),
sem augmentação. Dele sai a acurácia de referência $A^{\text{clean}}$, que
normaliza os números de robustez.

\subsection{Transformações não vistas}
\label{sec:protocol:unseen}
As corrupções da Tabela~\ref{tab:unseen} são aplicadas às imagens de teste em
$[0,1]$, depois da conversão para tensor e antes da normalização. A condição
``todas as transformações'' aplica as seis em sequência, na ordem da tabela.

\begin{table}[t]
  \centering
  \caption{Conjunto de transformações de teste. Nenhuma é aplicada em treino
  com esses parâmetros, mas algumas têm parentes próximos no regime forte (ver
  texto).}
  \label{tab:unseen}
  \begin{tabular}{lll}
    \toprule
    Transformação & Parâmetro & Valor \\
    \midrule
    Ruído gaussiano      & $\sigma$                & $0{,}08$ \\
    Desfoque gaussiano   & núcleo, $\sigma$        & $5$, $1{,}5$ \\
    Brilho               & fator                   & $1{,}6$ \\
    Contraste            & fator                   & $0{,}4$ \\
    Oclusão              & fração da área          & $0{,}25$ \\
    Rotação grande       & ângulo                  & $\theta \sim \mathcal{U}(-90^{\circ}, 90^{\circ})$ \\
    \bottomrule
  \end{tabular}
\end{table}

A separação entre treino e teste é completa para o regime padrão, mas parcial
para o forte. A oclusão cobre $25\%$ da imagem, dentro da faixa de área do
Random Erasing ($2\%$ a $33\%$, com probabilidade $0{,}5$). A diferença é que a
oclusão preenche o quadrado de preto, e o Random Erasing preenche o retângulo
com a cor média do conjunto. No D-Fire, o brilho $\times1{,}6$ também é
alcançável em treino: a variação de brilho vai até $\times1{,}3$ e o RandAugment
até $\times1{,}27$, e o produto chega a $\times1{,}65$. O contraste $\times0{,}4$
fica abaixo do mínimo de treino ($\times0{,}51$). Ruído e desfoque gaussianos
não têm equivalente no RandAugment; a operação de nitidez com fator
$0{,}73$ produz apenas um desfoque leve. Sob augmentação forte, portanto,
``não vista'' descreve a combinação e as magnitudes, não cada operação
isoladamente.

\subsection{Sonda de rotação}
\label{sec:protocol:rotation}
A rotação grande é o teste direto de H1 e também é avaliada isoladamente. Cada
imagem de teste recebe um único ângulo, sorteado de
$\mathcal{U}(-90^{\circ}, 90^{\circ})$, e o resultado é a acurácia agregada
sobre esses ângulos. Uma varredura determinística em $\theta$ produziria uma
curva por arquitetura; ela não foi feita, e a limitação é discutida na
Seção~\ref{sec:res:threats}.

No D-Fire, a rotação mede mais que equivariância. Chamas e plumas de fumaça
sobem, e esse prior de orientação é informação útil para o classificador.
Rotacionar a imagem quebra o prior, e a queda de acurácia mistura falta de
equivariância com perda dessa pista.

\subsection{Avaliação sob augmentação em teste}
Cada condição acima é avaliada duas vezes: sem augmentação em teste e com o
\emph{pipeline} de augmentação forte do conjunto aplicado às imagens de teste.
Trata-se de um deslocamento de distribuição adicional, e não de
\emph{test-time augmentation}: cada imagem passa pelo \emph{pipeline} uma
única vez, com operações sorteadas, e não há média de predições. As
corrupções não vistas entram depois do RandAugment e antes do Random Erasing.
São, portanto, seis condições por modelo treinado --- três transformações
$\times$ duas augmentações em teste ---, e 972 avaliações no total.

% ---------------------------------------------------------------------
\section{Métricas}
\label{sec:protocol:metrics}

\subsection{Qualidade da classificação}
A acurácia top-1 no conjunto de teste é a métrica principal. Ela é sempre lida
contra duas referências: o acaso ($10\%$ no CIFAR-10, $50\%$ no D-Fire) e a
linha de base da classe majoritária, que no D-Fire é de $53{,}4\%$ na partição
de teste. Um escore próximo dessas referências não é degradação graciosa; é um
modelo que deixou de discriminar.

Dois diagnósticos complementam a acurácia sob deslocamento, ambos definidos na
Seção~\ref{sec:res:collapse}: a AUROC da confiança do modelo como detector dos
próprios erros e a fração de predições concentradas na classe mais predita.
A taxa de falsos negativos da classe \textsc{fire/smoke}, relevante para um
detector de incêndio em operação, não é reportada: o objetivo aqui é comparar
arquiteturas, e não avaliar um sistema de detecção.

\subsection{Robustez}
Para uma corrupção $c$, a razão de retenção é
\begin{equation}
  R_{c} = \frac{A^{c}}{A^{\text{clean}}}.
  \label{eq:retention}
\end{equation}
A razão normaliza a robustez pela acurácia limpa, o que a torna tentadora para
comparar arquiteturas com desempenho em distribuição diferente. Ela só é
confiável quando as acurácias limpas são próximas: como $A^{c}$ tem piso no
acaso, um modelo fraco retém uma fração maior do pouco que tinha. Por isso a
degradação absoluta $A^{\text{clean}} - A^{c}$ e a habilidade retida $S_c$ da
Equação~\eqref{eq:skill}, que desconta o acaso, são reportadas ao lado.

\subsection{Eficiência de dados}
A estatística de eficiência de dados é a queda de acurácia limpa entre as
frações extremas:
\begin{equation}
  \Delta_{f} = A_{f=1{,}00} - A_{f=0{,}33}.
  \label{eq:data-eff}
\end{equation}
Com apenas três frações, a área sob a curva de acurácia contra fração não
acrescenta informação sobre $\Delta_f$ e não é usada.

\subsection{Custo computacional}
O custo de inferência está na Tabela~\ref{tab:model-complexity}. O custo de
treino, na Tabela~\ref{tab:cost}, vem dos registros por época: tempo de relógio
por época, horas por treino completo e pico de memória alocada pelo PyTorch na
GPU.

\begin{table}[t]
  \centering
  \caption{Custo de treino por arquitetura. Tempo por época e pico de memória
  são médias sobre os 27 treinos de cada célula; as horas por treino são média
  $\pm$ desvio padrão sobre esses treinos, e a dispersão vem sobretudo das três
  frações de dados.}
  \label{tab:cost}
  \begin{tabular}{llccc}
    \toprule
    Conjunto & Modelo & s/época & h/treino & Mem. (GB) \\
    \midrule
    \multirow{3}{*}{CIFAR-10}
      & Efficient-CapsNet & 15,6 & 0,4 $\pm$ 0,2 & 0,10 \\
      & ResNet-18         & 48,6 & 1,3 $\pm$ 0,5 & 0,27 \\
      & DeiT-Tiny         & 37,4 & 1,0 $\pm$ 0,4 & 0,23 \\
    \midrule
    \multirow{3}{*}{D-Fire}
      & Efficient-CapsNet & 49,4 & 1,4 $\pm$ 0,5 & 0,56 \\
      & ResNet-18         & 49,9 & 1,4 $\pm$ 0,5 & 0,63 \\
      & DeiT-Tiny         & 50,0 & 1,4 $\pm$ 0,5 & 0,54 \\
    \bottomrule
  \end{tabular}
\end{table}

No D-Fire, as três arquiteturas levam praticamente o mesmo tempo por época,
embora difiram em até quatro vezes no custo de inferência. O tempo desse
conjunto é dominado pela leitura e decodificação das imagens JPEG, feita por um
único processo auxiliar, e não pela rede. Os tempos do D-Fire, portanto, não
medem o custo das arquiteturas; os do CIFAR-10, sim.

% ---------------------------------------------------------------------
\section{Análise estatística}
\label{sec:protocol:stats}

Os resultados são reportados como média $\pm$ desvio padrão sobre as três
sementes. Com $n = 3$, o desvio padrão é uma estimativa fraca, e por isso a
variabilidade também é reportada como amplitude entre sementes
(Seção~\ref{sec:res:seeds}).

As comparações entre arquiteturas usam o teste $t$ pareado. O pareamento é por
semente, e, quando uma hipótese agrega várias condições, por semente e
condição: as três arquiteturas compartilham sementes, partições e fluxo de
augmentação (Seção~\ref{sec:protocol:repro}), e o pareamento remove essa fonte
comum de variação. Um teste não paramétrico foi descartado porque, com três
pares, o teste de Wilcoxon não pode produzir $p < 0{,}25$. O nível de
significância é $\alpha = 0{,}05$.

As hipóteses H1 e H2 são julgadas por seis comparações pareadas: retenção sob
rotação contra cada linha de base nos dois conjuntos, e $\Delta_f$ contra a
ResNet-18 nos dois conjuntos. Com a correção de Holm sobre essas seis
comparações, as diferenças contra o DeiT-Tiny continuam significativas
($p_{\text{Holm}} < 10^{-11}$). A desvantagem capsular sob rotação no D-Fire
passa de $p = 0{,}013$ para $p_{\text{Holm}} = 0{,}052$, e as demais ficam
entre $0{,}24$ e $0{,}26$. O veredito de nenhuma hipótese muda com a
correção, mas a afirmação de que a retenção capsular é \emph{menor} que a da
ResNet-18 no D-Fire deixa de ser estatisticamente significativa e deve ser
lida como tendência. As comparações célula a célula da
Seção~\ref{sec:res:wins} são descritivas e não recebem teste.

% ---------------------------------------------------------------------
\section{Reprodutibilidade}
\label{sec:protocol:repro}

Cada treino fixa a semente do PyTorch, que também semeia a GPU, e a do NumPy,
que controla a separação da validação e a subamostragem das frações. O cuDNN
roda em modo determinístico, com o \emph{benchmarking} desativado. O gerador do
PyTorch é semeado de novo imediatamente antes do laço de treino, depois da
construção do modelo, para que a inicialização de cada arquitetura, que consome
números aleatórios em quantidades diferentes, não altere o fluxo de
augmentação. Assim, para uma mesma semente, as três arquiteturas recebem os
mesmos lotes aumentados na mesma ordem, ao menos no início do treino.

Essa propriedade foi verificada reexecutando o procedimento de inicialização de
cada arquitetura, sem treinar, e comparando o \emph{hash} SHA-256 dos primeiros
lotes de treino.
% VERIFICAR: confirmar em paper_info/aug_determinism.csv que os hashes
% coincidem entre as três arquiteturas nos dois conjuntos.

A configuração YAML de cada execução é arquivada junto aos \emph{checkpoints},
e as métricas de treino e validação são gravadas por época em CSV. As
predições de teste são gravadas por imagem, o que permite recalcular todas as
tabelas do Capítulo~\ref{ch:results} sem reavaliar os modelos.

A tarefa de classificação usada no D-Fire não é a distribuída pelos autores. O
\emph{script} que converte as anotações YOLO nos rótulos binários da
Seção~\ref{sec:protocol:dfire} e a lista de arquivos com o rótulo e a partição
de cada imagem acompanham o código, junto com as configurações, em
\url{https://URL-DO-REPOSITORIO}.
% PREENCHER: endereço público do código.

% ---------------------------------------------------------------------
\section{Ameaças à validade}
\label{sec:protocol:threats}

\begin{description}

\item[Objetivos de treino distintos.]
O modelo de cápsulas usa perda de margem com reconstrução; as linhas de base usam
entropia cruzada. A diferença confunde a comparação arquitetural. Ela também não
é acessória: Gu, Tresp e Hu~\citar{gu2021notrobust} atribuem a essas perdas, e
não aos componentes capsulares, boa parte da robustez creditada às cápsulas.
Treinar as linhas de base com a mesma perda isolaria o efeito; essa variante
não foi executada.

\end{description}
